\documentclass[prd,twocolumn,floatfix,nofootinbib]{revtex4}
\usepackage{amssymb}
\usepackage{amsmath}
\usepackage{graphicx,color,dcolumn,booktabs,bm}
\usepackage{longtable,lscape}
\usepackage{txfonts}
\usepackage{overpic}
\usepackage{indentfirst}
\usepackage{cases}
\usepackage{multirow}
\usepackage{epstopdf}
\usepackage{ulem}
\usepackage[colorlinks,
            citecolor=blue,
            anchorcolor=red,
            menucolor=red,
            linkcolor=red,
            filecolor=red,
            runcolor=red,
            urlcolor=blue,
            frenchlinks=false]{hyperref}
%\newtheorem{theorem}{Theorem}
%\newtheorem{remark}[theorem]{Remark}
%\newtheorem{solution}[theorem]{Solution}
%\newtheorem{summary}[theorem]{Summary}
%\newenvironment{proof}[1][Proof]{\noindent\textbf{#1.} }{\ \rule{0.5em}{0.5em}}
%\usepackage{tikz}
\setcounter{MaxMatrixCols}{10}
\def\lrpartial{\buildrel\leftrightarrow\over\partial}
\graphicspath{{Figures/}} %
\allowdisplaybreaks

\linespread{1.1}
\begin{document}
\title{Ground-state baryons and tetraquarks with one and two heavy quarks: Masses, magnetic moments}

\author{Wen-Xuan Zhang$^{1}$}
\author{Hao Xu$^{1,2}$}
\author{Duojie Jia$^{1,2,3,4}$
\footnote{Corresponding author}}
\email{jiadj@nwnu.edu.cn}

\affiliation{ $^1$Institute of Theoretical Physics, College of Physics and
Electronic Engineering, Northwest Normal University, Lanzhou 730070, China \\
$^2$Lanzhou Center for Theoretical Physics, Lanzhou University, Lanzhou, 730000,China \\
$^3$College of Physics and Electronic Information Engineering,
Qinghai Normal University, Xining 810000, China \\
$^4$Joint Research Center for Physics, Lanzhou University
and Qinghai Normal University, Xining 810000,China \\ }


\begin{abstract}
%Recently, LHCb Collaboration observed a very prominate resonances $X_{0}(2900)$ in $B^{+}\to D^{+}D^{-}K^{+}$ decays, which implies open-heavy-flavor exotic hadron $ud\bar{s}\bar{c}$.
%Motivated by this, we estimate masses and magnetic moments of ground-state baryons with one or two heavy quarks, in an unified description of MIT bag model for all light and heavy hadrons.
%For heavy sectors, we use masses of heavy mesons with $J^{P}=1^{-}$, to determine the effective masses of charm and bottom quarks, as well as binding energies which are assumed to exist among them and strange quarks.
%The mass and magnetic moments of $\Xi_{cc}^{\ast}$ are predicted to be $M\left(\Xi_{cc}^{\ast}\right)=3.714\,{\rm GeV}$, $\mu\left(\Xi_{cc}^{\ast ++}\right)=0.94\mu_{p}$ and $\mu\left(\Xi_{cc}^{\ast +}\right)=0.06\mu_{p}$.
%The mass of isosinglet $ud\bar{s}\bar{c}$ with $J^{P}=0^{+}$ are predicted to be $M\left(ud\bar{s}\bar{c}\right)=2.934\,{\rm GeV}$.
In this paper, we estimate 

{\normalsize PACS number(s):12.39Jh, 12.40.Yx, 12.40.Nn}

{\normalsize Key Words: Heavy tetraquark, Spectroscopy, Quantum number, Magnetic moment}
\end{abstract}
\maketitle
\date{\today}


\maketitle\section{Introduction}
This paper will be organized as following: In Sect.~\ref{Method}, we introduce some basic formulae of MIT bag model with which one is familiar in other people's work and of CMI model to determine spin-dependent interactions for considered hadrons. 
In Sect.~\ref{Parameters}, we perform a series of numerical calculations to find out the optimal set of parameters which is good for fitting hadrons' spectra. 
In Sect.~\ref{Hadrons}, these parameters have already been prepared to calculate the mass spectrum, magnetic moments and charge radius for doubly heavy, triply heavy baryons and tetraquarks. 
We shall discuss these systems in separated subsections for details. 
Then Sect.~\ref{Conclusions} provides analysis and comparations about hadron's data with the one in other people's work. 
And in the same section, we turn back to the physics to bring out more useful ideas.


\section{Method for MIT bag model and CMI}\label{Method}
\subsection{Mass Formula}
In MIT bag model, the mass formula of hadron is
\begin{equation}\label{MBm}
    M\left(R\right)=\sum_{i=n,s,c,b} n_{i}\omega_{i}+\frac{4}{3}\pi R^{3}B-\frac{Z_{0}}{R}+\Delta M
\end{equation}
\begin{equation}\label{freq}
    \omega_{i}=\left(m_{i}^{2}+\frac{x_{i}^{2}}{R^2}\right)^{1/2}
\end{equation}
where $m_{i}$ is mass of quark or antiquark for $n$ (denotes $u$, $d$), $s$, $c$, $b$ flavor, $x_{i}=x_{i}(mR)$ is dimensionless dynamical parameters, 
$R$ is bag radius, $B$ is bag constant and $Z_{0}$ is zero-point-energy (ZPE) constant in its nature. 
$\Delta M$ is interaction energy which we address in this work.
We choose $m_{i}$, $B$, $Z_{0}$ as free parameters and treat $x_{i}$, $R$ as intermediate quantities satisfying a transcendental equation
\begin{equation}\label{transc}
    \tan x_{i}=\frac{x_{i}}{1-m_{i}R-\left(m_{i}^{2}R^{2}+x_{i}^{2}\right)^{1/2}}
\end{equation}

The first part of hadron mass in Eq.~(\ref{MBm}) is sum of relativistical energy of quarks contained in the bag.
With the interaction energy $\Delta M$ including Coulomb-like and spin-dependent potentials, it contributes to ordinary quark model.
But MIT bag model offers more physical pictures about QCD, 
that the second part is volume energy in which constant $B$ standing for color confinement indicates difference of energy density between perturbation and non-perturbation QCD vacuum.
And the third part is ZPE produced by Casimir effect in a sphere bag.

The mass formula (\ref{MBm}) is a function of bag radius $R$, and most terms in it are also $R$-dependent.
Because it’s believed that a sphere bag surface can be balanced by Dirac pressure inside and vacuum condensation outside. 
For the stability, bag energy will locate at minimum with respect to bag radius $R$. 
That means we need only a few free parameters to calculate all the lowest-lying hadronic states, from mesons to tetraquarks, by variational method.
And we would rather consider recoil of heavy quark, than taking replacement of the first term in Eq.~(\ref{MBm}) by fixed constituent mass.

The interaction energy $\Delta M$ in Eq.~(\ref{MBm}) contains Coulomb-like and spin-dependent potentials.
In Refs.~[], Karliner has introduced binding effect between two relatively heavy quarks, in studying ordinary quark model.
This effect is thought of as short-disctance chromoelectric interactions between c and s, c and c, b and s, b and b, b and c quarks respectively, 
which form strong bindings namely binding energy. In this work, we will have five bindings $B_{cs}$, $B_{cc}$, $B_{bs}$, $B_{bb}$, $B_{bc}$ for color $\boldsymbol{\bar{3}}$, 
and scale them by color factors to other color configurations.

The second part of $\Delta M$ is chromomagnetic interaction
\begin{equation}\label{CMI}
    H_{CMI}=-\sum_{i<j} \left(\boldsymbol{\lambda_{i}}\cdot \boldsymbol{\lambda_{j}}\right) \left(\boldsymbol{\sigma_{i}}\cdot \boldsymbol{\sigma_{j}}\right)C_{ij}
\end{equation}
where $i$ and $j$ denote quarks or antiquarks, $\boldsymbol{\lambda}$ and $\boldsymbol{\sigma}$ are Gell-Mann and Pauli matrices respectively. 
$C_{ij}$ is dynamical parameter which can be derived by MIT bag model as
\begin{equation}\label{Cij}
    C_{ij}=3\frac{\alpha_{s}\left(R\right)}{R^{3}}\bar{\mu}_{i}\bar{\mu}_{j}I_{ij}
\end{equation}
\begin{equation}\label{muBari}
   \bar{\mu}_{i}=\frac{R}{6} \frac{4\alpha_{i}+2\lambda_{i}-3}{2\alpha_{i}\left(\alpha_{i}-1\right)+\lambda_{i}}
\end{equation}
\begin{equation}\label{Iij}
    I_{ij}=1+2\int_{0}^{R} \frac{{\rm d}r}{r^{4}} \bar{\mu}_{i}\bar{\mu}_{j}=
    1+\text{rational term about } x_{i} \text{, } x_{j}
\end{equation}
where $\alpha_{i}\equiv\omega_{i}R$, $\lambda_{i}\equiv m_{i}R$, $\alpha_{s}\left(R\right)$ is running strong coupling constant, $\bar{\mu}_{i}$ is reduced magnetic moment without quark charge,
and $I_{ij}$ denoting improper integral has a complicated rational expression one can find in Ref.~[]. 
In Sect.~\ref{Parameters}, $\alpha_{s}\left(R\right)$ will be discussed and chose for fitting spectrum well.


\subsection{Chromomagnetic Structure}
The chromomagnetic structure for a hadronic state is required, if one needs to determine the expression of Eq.~(\ref{CMI}) after identifying this state with certian $J^{P}$ quantum number.
In this section, we will mainly introduce the flavor-spin-color wave functions for meson, baryon and tetraquark systems.

Mesons have color singlet and two spin states
\begin{equation}\label{colorM}
    \phi^{M}=\left| q_{1}\bar{q}_{2} \right\rangle
\end{equation}
\begin{equation}\label{spinM}
    \chi_{1}^{M}={\left| q_{1}\bar{q}_{2} \right\rangle}_{1} \quad \chi_{2}^{M}={\left| q_{1}\bar{q}_{2} \right\rangle}_{0}
\end{equation}
where subscript outside the bracket denotes quantum number $J$. The corresponding combinations give
\begin{equation}\label{colorspinM}
    \phi^{M}\chi_{1}^{M}={\left| q_{1}\bar{q}_{2} \right\rangle}_{1} \quad \phi^{M}\chi_{2}^{M}={\left| q_{1}\bar{q}_{2} \right\rangle}_{0}
\end{equation}

Baryons have color $\boldsymbol{\bar{3}}$ and three spin states
\begin{equation}\label{colorB}
    \phi^{B}=\left| {\left(q_{1}q_{2}\right)}^{\bar{3}} q_{3} \right\rangle
\end{equation}
\begin{equation}\label{spinB}
    \begin{aligned}
        \chi_{1}^{B}={\left| {\left(q_{1}q_{2}\right)}_{1} q_{3} \right\rangle}_{3/2} \\ 
        \chi_{2}^{B}={\left| {\left(q_{1}q_{2}\right)}_{1} q_{3} \right\rangle}_{1/2} \\
        \chi_{3}^{B}={\left| {\left(q_{1}q_{2}\right)}_{0} q_{3} \right\rangle}_{1/2}
    \end{aligned}
\end{equation}
where $\left(q_{1}q_{2}\right)$ is diquark with spin number 1 or 0.

We now move to the flavor symmetry.
If the flavor wave function of $\left(q_{1}q_{2}\right)$ is symmetric, when isospin $I=1$ or flavors are identical, we choose a symbol $\delta_{12}^{S}=1$. 
Otherwise, if the flavor is asymmetric when $I=0$, there will be $\delta_{12}^{A}=1$. 
And if $q_{q}$ and $q_{2}$ carry different flavors beyond SU(2) symmetry, we believe the wave function can be either symmetric or asymmetric namely $\delta_{12}^{S}=\delta_{12}^{A}=1$. 
Consequently, we obtain three flavor-spin-color wave functions according to Pauli principle writing in
\begin{equation}\label{colorspinB}
    \begin{aligned}
        \phi^{B}\chi_{1}^{B}={\left| {\left(q_{1}q_{2}\right)}_{1}^{\bar{3}} q_{3} \right\rangle}_{3/2} \delta_{12}^{S} \\
        \phi^{B}\chi_{2}^{B}={\left| {\left(q_{1}q_{2}\right)}_{1}^{\bar{3}} q_{3} \right\rangle}_{1/2} \delta_{12}^{S} \\
        \phi^{B}\chi_{3}^{B}={\left| {\left(q_{1}q_{2}\right)}_{0}^{\bar{3}} q_{3} \right\rangle}_{1/2} \delta_{12}^{A} 
    \end{aligned}
\end{equation}

In this work, we also study chromomagnetic mixing. 
For example, $\Xi_{bc}$ and $\Xi_{bc}^{\prime}$ with flavor structure $(bc)s$ choosing $\delta_{12}^{S}=\delta_{12}^{A}=1$ can't distinguish $\phi^{B}\chi_{2}^{B}$ and $\phi^{B}\chi_{3}^{B}$ states by $J^{P}$ quantum number, 
thus the physical state must be mixing of $\phi^{B}\chi_{2}^{B}$ and $\phi^{B}\chi_{3}^{B}$.
We will make full discussion about states mixing in Sect.~\ref{Parameters}.

Finally, tetraquarks have color $\boldsymbol{6}\otimes\boldsymbol{\bar{6}}$ and $\boldsymbol{\bar{3}}\otimes\boldsymbol{3}$ and six spin states
\begin{equation}\label{colorT}
    \phi_{1}^{T}=\left| {\left(q_{1}q_{2}\right)}^{6} {\left(\bar{q}_{3}\bar{q}_{4}\right)}^{\bar{6}} \right\rangle \quad
    \phi_{2}^{T}=\left| {\left(q_{1}q_{2}\right)}^{\bar{3}} {\left(\bar{q}_{3}\bar{q}_{4}\right)}^{3} \right\rangle
\end{equation}
\begin{equation}\label{spinT}
    \begin{aligned}
        \chi_{1}^{T}={\left| {\left(q_{1}q_{2}\right)}_{1} {\left(\bar{q}_{3}\bar{q}_{4}\right)}_{1} \right\rangle}_{2} \quad
        \chi_{2}^{T}={\left| {\left(q_{1}q_{2}\right)}_{1} {\left(\bar{q}_{3}\bar{q}_{4}\right)}_{1} \right\rangle}_{1} \\
        \chi_{3}^{T}={\left| {\left(q_{1}q_{2}\right)}_{1} {\left(\bar{q}_{3}\bar{q}_{4}\right)}_{1} \right\rangle}_{0} \quad
        \chi_{4}^{T}={\left| {\left(q_{1}q_{2}\right)}_{1} {\left(\bar{q}_{3}\bar{q}_{4}\right)}_{0} \right\rangle}_{1} \\
        \chi_{5}^{T}={\left| {\left(q_{1}q_{2}\right)}_{0} {\left(\bar{q}_{3}\bar{q}_{4}\right)}_{1} \right\rangle}_{1} \quad
        \chi_{6}^{T}={\left| {\left(q_{1}q_{2}\right)}_{0} {\left(\bar{q}_{3}\bar{q}_{4}\right)}_{0} \right\rangle}_{0}
    \end{aligned}
\end{equation}
with twelve wave functions
\begin{equation}\label{colorspinT}
    \begin{aligned}
        \phi_{1}^{T}\chi_{1}^{T}={\left| {\left(q_{1}q_{2}\right)}_{1}^{6} {\left(\bar{q}_{3}\bar{q}_{4}\right)}_{1}^{\bar{6}} \right\rangle}_{2} \delta_{12}^{A}\delta_{34}^{A} \\
        \phi_{2}^{T}\chi_{1}^{T}={\left| {\left(q_{1}q_{2}\right)}_{1}^{\bar{3}} {\left(\bar{q}_{3}\bar{q}_{4}\right)}_{1}^{3} \right\rangle}_{2} \delta_{12}^{S}\delta_{34}^{S} \\
        \phi_{1}^{T}\chi_{2}^{T}={\left| {\left(q_{1}q_{2}\right)}_{1}^{6} {\left(\bar{q}_{3}\bar{q}_{4}\right)}_{1}^{\bar{6}} \right\rangle}_{1} \delta_{12}^{A}\delta_{34}^{A} \\
        \phi_{2}^{T}\chi_{2}^{T}={\left| {\left(q_{1}q_{2}\right)}_{1}^{\bar{3}} {\left(\bar{q}_{3}\bar{q}_{4}\right)}_{1}^{3} \right\rangle}_{1} \delta_{12}^{S}\delta_{34}^{S} \\
        \phi_{1}^{T}\chi_{3}^{T}={\left| {\left(q_{1}q_{2}\right)}_{1}^{6} {\left(\bar{q}_{3}\bar{q}_{4}\right)}_{1}^{\bar{6}} \right\rangle}_{0} \delta_{12}^{A}\delta_{34}^{A} \\
        \phi_{2}^{T}\chi_{3}^{T}={\left| {\left(q_{1}q_{2}\right)}_{1}^{\bar{3}} {\left(\bar{q}_{3}\bar{q}_{4}\right)}_{1}^{3} \right\rangle}_{0} \delta_{12}^{S}\delta_{34}^{S} \\
        \phi_{1}^{T}\chi_{4}^{T}={\left| {\left(q_{1}q_{2}\right)}_{1}^{6} {\left(\bar{q}_{3}\bar{q}_{4}\right)}_{0}^{\bar{6}} \right\rangle}_{1} \delta_{12}^{A}\delta_{34}^{S} \\
        \phi_{2}^{T}\chi_{4}^{T}={\left| {\left(q_{1}q_{2}\right)}_{1}^{\bar{3}} {\left(\bar{q}_{3}\bar{q}_{4}\right)}_{0}^{3} \right\rangle}_{1} \delta_{12}^{S}\delta_{34}^{A} \\
        \phi_{1}^{T}\chi_{5}^{T}={\left| {\left(q_{1}q_{2}\right)}_{0}^{6} {\left(\bar{q}_{3}\bar{q}_{4}\right)}_{1}^{\bar{6}} \right\rangle}_{1} \delta_{12}^{S}\delta_{34}^{A} \\
        \phi_{2}^{T}\chi_{5}^{T}={\left| {\left(q_{1}q_{2}\right)}_{0}^{\bar{3}} {\left(\bar{q}_{3}\bar{q}_{4}\right)}_{1}^{3} \right\rangle}_{1} \delta_{12}^{A}\delta_{34}^{S} \\
        \phi_{1}^{T}\chi_{6}^{T}={\left| {\left(q_{1}q_{2}\right)}_{0}^{6} {\left(\bar{q}_{3}\bar{q}_{4}\right)}_{0}^{\bar{6}} \right\rangle}_{0} \delta_{12}^{S}\delta_{34}^{S} \\
        \phi_{2}^{T}\chi_{6}^{T}={\left| {\left(q_{1}q_{2}\right)}_{0}^{\bar{3}} {\left(\bar{q}_{3}\bar{q}_{4}\right)}_{0}^{3} \right\rangle}_{0} \delta_{12}^{A}\delta_{34}^{A}
    \end{aligned}
\end{equation}

Appendix A and B present all the precise color, spin wave functions. Therefore, the color and spin factors in Eq.~(\ref{CMI}) can be calculated by using
\begin{equation}\label{colorfc}
    {\left\langle \boldsymbol{\lambda_{i}}\cdot \boldsymbol{\lambda_{j}} \right\rangle}_{nm}=
    \sum_{\alpha=1}^{8} Tr\left(c_{in}^{\dagger}\lambda^{\alpha}c_{im}\right) Tr\left(c_{jn}^{\dagger}\lambda^{\alpha}c_{jm}\right)
\end{equation}
\begin{equation}\label{spinfc}
    {\left\langle \boldsymbol{\sigma_{i}}\cdot \boldsymbol{\sigma_{j}} \right\rangle}_{xy}=
    \sum_{\alpha=1}^{3} Tr\left(\chi_{ix}^{\dagger}\sigma^{\alpha}\chi_{iy}\right) Tr\left(\chi_{jx}^{\dagger}\sigma^{\alpha}\chi_{jy}\right)
\end{equation}
where $n$, $m$ and $x$, $y$ indicate the specific color and spin states respectively, $i$, $j$ are indexes of quarks or antiquarks and $c$, $\chi$ are basis vectors for color and spin space.
Several non-mixing hadronic states with CMI are listed in Table \ref{tab:CMI}.
We are now able to construct the whole mass formula (\ref{MBm}) and evaluate mass spectrum and dynamical quantities $x_{i}$, $R$.

\renewcommand{\tabcolsep}{0.49cm}
\renewcommand{\arraystretch}{1.5}
\begin{table*}[!htbp]
    \caption{CMI for several non-mixing hadrons with chromomagnetic structures. $C_{ij}$ follows Eq.~(\ref{Cij}) with subscripts which do not differ quark and antiquark.}
    \label{tab:CMI}
    \begin{tabular}{cccccc}
        \hline\hline
        {\rm State}     &{\rm Structure}            &{\rm CMI}      &{\rm State}        &{\rm Structure}            &{\rm CMI} \\ \hline
        $\pi$           &$\phi^{M}\chi_{2}^{M}$     &$-16C_{nn}$    &$\omega$           &$\phi^{M}\chi_{1}^{M}$     &$\frac{16}{3}C_{nn}$ \\ 
        $K$             &$\phi^{M}\chi_{2}^{M}$     &$-16C_{sn}$    &$K^{\ast}$         &$\phi^{M}\chi_{1}^{M}$     &$\frac{16}{3}C_{sn}$ \\ 
                        &                           &               &$\phi$             &$\phi^{M}\chi_{1}^{M}$     &$\frac{16}{3}C_{ss}$ \\ 
        $D$             &$\phi^{M}\chi_{2}^{M}$     &$-16C_{cn}$    &$D^{\ast}$         &$\phi^{M}\chi_{1}^{M}$     &$\frac{16}{3}C_{cn}$ \\ 
        $D_{s}$         &$\phi^{M}\chi_{2}^{M}$     &$-16C_{cs}$    &$D_{s}^{\ast}$     &$\phi^{M}\chi_{1}^{M}$     &$\frac{16}{3}C_{cs}$ \\ 
        $\eta_{c}$      &$\phi^{M}\chi_{2}^{M}$     &$-16C_{cc}$    &$J/\psi$           &$\phi^{M}\chi_{1}^{M}$     &$\frac{16}{3}C_{cc}$ \\ 
        
        $N$             &$\phi^{B}\chi_{3}^{B}$     &$-8C_{nn}$                                 &$\Delta$               &$\phi^{B}\chi_{1}^{B}$     &$8C_{nn}$ \\ 
        $\Lambda$       &$\phi^{B}\chi_{3}^{B}$     &$-8C_{nn}$                                 &                       &                           & \\ 
        $\Sigma$        &$\phi^{B}\chi_{2}^{B}$     &$\frac{8}{3}C_{nn}-\frac{32}{3}C_{sn}$     &$\Sigma^{\ast}$        &$\phi^{B}\chi_{1}^{B}$     &$\frac{8}{3}C_{nn}+\frac{16}{3}C_{sn}$ \\ 
        $\Lambda_{c}$   &$\phi^{B}\chi_{3}^{B}$     &$-8C_{nn}$                                 &                       &                           & \\ 
        $\Sigma_{c}$    &$\phi^{B}\chi_{2}^{B}$     &$\frac{8}{3}C_{nn}-\frac{32}{3}C_{cn}$     &$\Sigma_{c}^{\ast}$    &$\phi^{B}\chi_{1}^{B}$     &$\frac{8}{3}C_{nn}+\frac{16}{3}C_{cn}$ \\ 
        $\Xi$           &$\phi^{B}\chi_{2}^{B}$     &$\frac{8}{3}C_{ss}-\frac{32}{3}C_{sn}$     &$\Xi^{\ast}$           &$\phi^{B}\chi_{1}^{B}$     &$\frac{8}{3}C_{ss}+\frac{16}{3}C_{sn}$ \\ 
        $\Xi_{c}$       &$\phi^{B}\chi_{3}^{B}$     &$-8C_{sn}$                                 &$\Xi_{c}^{\ast}$       &$\phi^{B}\chi_{1}^{B}$     &$\frac{8}{3}C_{sn}+\frac{8}{3}C_{cn}+\frac{8}{3}C_{cs}$ \\ 
                        &                           &                                           &$\Omega$               &$\phi^{B}\chi_{1}^{B}$     &$8C_{ss}$ \\ 
        
        ${\left(nn\bar{c}\bar{c}\right)}^{I=1}$     &$\phi_{2}^{T}\chi_{2}^{T}$     &$\frac{8}{3}C_{nn}-\frac{16}{3}C_{cn}+\frac{8}{3}C_{cc}$
        &${\left(nn\bar{c}\bar{c}\right)}^{I=1}$    &$\phi_{2}^{T}\chi_{1}^{T}$     &$\frac{8}{3}C_{nn}+\frac{16}{3}C_{cn}+\frac{8}{3}C_{cc}$ \\ 
        ${\left(nn\bar{s}\bar{c}\right)}^{I=0}$     &$\phi_{1}^{T}\chi_{1}^{T}$     &$-\frac{4}{3}C_{nn}+\frac{20}{3}C_{sn}+\frac{20}{3}C_{cn}-\frac{4}{3}C_{cs}$
        &${\left(nn\bar{s}\bar{c}\right)}^{I=1}$    &$\phi_{2}^{T}\chi_{1}^{T}$     &$\frac{8}{3}C_{nn}+\frac{8}{3}C_{sn}+\frac{8}{3}C_{cn}+\frac{8}{3}C_{cs}$ \\ 
        \hline\hline
    \end{tabular}
\end{table*}


\subsection{Hadronic Properties}
After obtaining the quantities $x_{i}$, $R$ for a hadronic state, we subsequently calculate charge radius and magnetic moment,
which depend on dynamical information of quarks.
We first use the following mothed to evaluate contribution of quarks on charge radius
\begin{equation}\label{rEi}
    \begin{aligned}
        {\left\langle r_{E}^{2} \right\rangle}_{i}&=
        Q_{i}R^{2} \frac{\alpha_{i} \left[2x_{i}^{2}\left(\alpha_{i}-1\right)+4\alpha_{i}+2\lambda_{i}-3\right]}
        {3x_{i}^{2} \left[2\alpha_{i}\left(\alpha_{i}-1\right)+\lambda_{i}\right]} \\ 
        &-Q_{i}R^{2} \frac{\lambda_{i} \left[4\alpha_{i}+2\lambda_{i}-2x_{i}^{2}-3\right]}
        {2x_{i}^{2} \left[2\alpha_{i}\left(\alpha_{i}-1\right)+\lambda_{i}\right]}
    \end{aligned}
\end{equation}
where $Q_{i}$ is charge of quark or antiquark. 
Since the spin-color wave functions cannot apply to charge radius, one can directly sum up Eq.~(\ref{rEi}) by
\begin{equation}\label{rEsum}
    r_{E}={\left| \sum\nolimits_{q} {\left\langle r_{E}^{2} \right\rangle}_{q} \right|}^{1/2}
\end{equation}
to give out the charge radius of hadronic state.
Eq.~(\ref{rEsum}) is also correct for mixing systems which have the same quark constituent.

We then use the following equation to calculate magnetic moment
\begin{equation}\label{mui}
    \mu_{i}=Q_{i}\bar{\mu}_{i}=Q_{i}\frac{R}{6} \frac{4\alpha_{i}+2\lambda_{i}-3}{2\alpha_{i}\left(\alpha_{i}-1\right)+\lambda_{i}}
\end{equation}
\begin{equation}\label{musum}
    \mu=\left\langle \psi_{spin} \left| \sum\nolimits_{i}g_{i}\mu_{i}S_{iz} \right| \psi_{spin} \right\rangle
\end{equation}
where $S_{iz}$ is the third component of spin for individual quark or antiquark, and $g_{i}=2$ is a Lander factor. 
Moreover, if mixing states are took into account whose total spin wave function reads
\begin{equation}
    \left| \psi_{spin} \right\rangle=C_{1}\chi_{1}+C_{2}\chi_{2}
\end{equation}
one can express magnetic moment as
\begin{equation}\label{mixingmu}
    \mu=C_{1}^{2}\mu\left(\chi_{1}\right)+C_{2}^{2}\mu\left(\chi_{2}\right)+2C_{1}C_{2}\mu^{\rm tr}\left(\chi_{1}, \chi_{2}\right)
\end{equation}
where $\mu^{\rm tr}$ is transition moment studied in [ref] and $C_{1}$, $C_{2}$ form eigenvector for mixing states.
According to spin wave functions in Appendix B, we derive magnetic moments for all the spin states and mixing systems involved in this work, and list them in Table \ref{tab:musum}.

We notice that, the transition moments in mixing systems are not symmetric by swaping $q_{1}$, $q_{2}$ or $q_{3}$, $q_{4}$ in flavor wave functions respectively.
While isospin shall not be considered in constructing mass formula, it accually distinguishs magnetic moments, between $\left(ud\right)s$, $\left(du\right)s$ states for example,
by the isospin wave function
\begin{equation}
    {\left| 1,0\right\rangle}_{I}=\frac{1}{\sqrt{2}} \left(ud+du\right) \quad
    {\left| 0,0\right\rangle}_{I}=\frac{1}{\sqrt{2}} \left(ud-du\right)
\end{equation}
One must take average of these two states to evaluate magnetic moment physically, but consequently the transition moments vanish.
So there is no transition moment between two mixing state, and we will keep only the terms with $C_{i}^{2}$ in Eq.~(\ref{mixingmu}).

At the end of this section, we further mention that, for a state with precise flavor-spin-color structure and certian CMI, it has a unique set of $x_{i}$, $R$ comparing with others, even those states carring the same $J^{P}$ quantum number.
It restricts us to determine mass formula, for mixinig systems indicating analytical eigenvalues of CMI matrix, before performing variation to it.
One may keep this in mind when calculating mass spectrum and hadronic properties of mixing states especially tetraquarks.

\renewcommand{\tabcolsep}{0.7cm}
\renewcommand{\arraystretch}{1.5}
\begin{table*}[!htbp]
    \caption{Sum rule of magnetic moment for a given spin state or a mixing system.}
    \label{tab:musum}
    \begin{tabular}{lcc}
        \hline\hline
        {\rm Hadron}        &$\psi_{spin}$      &$\mu$ \\ \hline
        $\left(q_{1}{\bar{q}}_{2}\right)$
                            &$\chi_{1}^{M}$     &$\mu_{1}+\mu_{2}$ \\ 
                            &$\chi_{2}^{M}$     &0 \\ 
        $\left(q_{1}q_{2}\right)q_{3}$
                            &$\chi_{1}^{B}$     &$\mu_{1}+\mu_{2}+\mu_{3}$ \\ 
                            &$\chi_{2}^{B}$     &$\frac{1}{3}\left(2\mu_{1}+2\mu_{2}-\mu_{3}\right)$ \\ 
                            &$\chi_{3}^{B}$     &$\mu_{3}$ \\ 
                            &$C_{1}\chi_{2}^{B}+C_{2}\chi_{3}^{B}$      
                            &$C_{1}^{2}\mu\left(\chi_{2}^{B}\right) + C_{2}^{2}\mu\left(\chi_{3}^{B}\right) + \frac{2}{\sqrt{3}}C_{1}C_{2}\left(\mu_{2}-\mu_{1}\right)$ \\ 
        $\left(q_{1}q_{2}\right)\left({\bar{q}}_{3}{\bar{q}}_{4}\right)$
                            &$\chi_{1}^{T}$     &$\mu_{1}+\mu_{2}+\mu_{3}+\mu_{4}$ \\ 
                            &$\chi_{2}^{T}$     &$\frac{1}{2}\left(\mu_{1}+\mu_{2}+\mu_{3}+\mu_{4}\right)$ \\ 
                            &$\chi_{3}^{T}$     &0 \\ 
                            &$\chi_{4}^{T}$     &$\mu_{1}+\mu_{2}$ \\ 
                            &$\chi_{5}^{T}$     &$\mu_{3}+\mu_{4}$ \\ 
                            &$\chi_{6}^{T}$     &0 \\ 
                            &$C_{1}\chi_{3}^{T}+C_{2}\chi_{6}^{T}$      &0 \\ 
                            &$C_{1}\chi_{5}^{T}+C_{2}\chi_{4}^{T}$
                            &$C_{1}^{2}\mu\left(\chi_{5}^{T}\right) + C_{2}^{2}\mu\left(\chi_{4}^{T}\right)$ \\ 
                            &$C_{1}\chi_{2}^{T}+C_{2}\chi_{4}^{T}+C_{3}\chi_{5}^{T}$
                            &$C_{1}^{2}\mu\left(\chi_{2}^{T}\right) + C_{2}^{2}\mu\left(\chi_{4}^{T}\right) + C_{3}^{2}\mu\left(\chi_{5}^{T}\right)
                            + \sqrt{2}C_{1}C_{2}\left(\mu_{3}-\mu_{4}\right) + \sqrt{2}C_{1}C_{3}\left(\mu_{2}-\mu_{1}\right)$ \\ 
                            &$C_{1}\chi_{2}^{T}+C_{2}\chi_{5}^{T}+C_{3}\chi_{4}^{T}$
                            &$C_{1}^{2}\mu\left(\chi_{2}^{T}\right) + C_{2}^{2}\mu\left(\chi_{5}^{T}\right) + C_{3}^{2}\mu\left(\chi_{4}^{T}\right)
                            + \sqrt{2}C_{1}C_{2}\left(\mu_{2}-\mu_{1}\right) + \sqrt{2}C_{1}C_{3}\left(\mu_{3}-\mu_{4}\right)$ \\ 
        \hline\hline
    \end{tabular}
\end{table*}


\section{Determination of Parameters}\label{Parameters}
MIT bag model only discusses a few free parameters such as quark mass $m_{i}$, bag constant $B$, coefficient $Z_{0}$ of ZPE and strong coupling constant $\alpha_{s}$.
We quote the original work of this model, where the parameters are determined so that light hadrons $N$, $\Delta$, $\omega$ and $\Omega$ fit in spectrum, that suggests 
\begin{equation}\label{originalparas}
    \begin{Bmatrix}
        m_{n}=0\,{\rm GeV} & m_{s}=0.279\,{\rm GeV} \\
        Z_{0}=1.84 & B^{1/4}=0.145\,{\rm GeV} & \alpha_{s}=0.55
    \end{Bmatrix}
\end{equation}

Some people would like to refine these conclusions to reproduce spectrum of heavy hadrons with charm and bottom quarks in Refs.~[].
Our considerations are: the parameters of MIT bag model are hard to be determined because tuning them always fails to cover all the confirmed states;
parameters describing physical properties in the same model shouldn't vary a lot.
Thus we will keep Eq.~(\ref{originalparas}) with little change in our work to calculate all the lowest-lying states, as this model is relatively general.

As noted above, free parameters are hard to reproduce all the states.
$B$ and $Z_{0}$ that largely lift up or lower down the whole spectrum cannot reduce all the errors simultaneously, 
while light baryons with strange flavor and $K^{\ast}$ in Ref.~[] will not get close to zero-error together by changing $m_{s}$.
The only thing we believe that provides flexible transition is running strong coupling constant $\alpha_{s}$.

Inspired by QCD renormalization, $\alpha_{s}$ takes the semiempirical form in MIT bag model as 
\begin{equation}\label{alphaS}
    \alpha_{s}(R)=\frac{2\pi n}{9} \frac{1}{ln\left[\gamma+{\left(R\Lambda_{QCD}\right)}^{-n}\right]}
\end{equation}
where exponent $n$ can be either 1 or 2, while 1 is prefered for suppressing the value of $\alpha_{s}$,
$\Lambda_{QCD}$ denotes QCD scalar taking value from $0.2\,{\rm GeV}$ to $0.5\,{\rm GeV}$,
and $\gamma$ is just a phenomenological variate meaning to avoid divergence.
Several expressions have been used in different works [], most of which only change $\gamma$, $\Lambda_{QCD}$ and $n$.
But there is a new form
\begin{equation}\label{alphaS-2004}
    \alpha_{s}(R)=\frac{0.470}{ln\left[1+{\left(0.420R\right)}^{-1}\right]}
\end{equation}
according to Da-Heng He's work [], where $0.470$ is a coefficient, and $\Lambda_{QCD}=0.420\,{\rm GeV}$.

The problem emerges if we put Eq.~(\ref{alphaS-2004}) into Eq.~(\ref{Cij}), the errors of light hadrons will be enlarged significantly, since $\alpha_{s}=1.21$ when $R=5\,{\rm GeV}^{-1}$.
It's too large comparing with Eq.~(\ref{originalparas}), and strong coupling constant is quite smaller in pQCD actually.
For physical concerning, a renormalization scalar $\mu=1\,{\rm GeV}$ will fit $\alpha_{s}=0.55$ by zero-order approximation of QCD, getting $\Lambda_{QCD}=0.281\,{\rm GeV}$. 
Then we set $R=5\,{\rm GeV}^{-1}$ for aligning proton in MIT bag model, and fit 0.55 again by the form like Eq.~(\ref{alphaS-2004}) to come up with our conclusion
\begin{equation}\label{alphaS-mine}
    \alpha_{s}(R)=\frac{0.296}{ln\left[1+{\left(0.281R\right)}^{-1}\right]}
\end{equation}

In Fig.~\ref{fig:alphaS}, there are two more running behaviors comparing to Eq.~(\ref{alphaS-2004}) and Eq.~(\ref{alphaS-mine}),
in which only solid line provides effective variation with lower value at $R=5\,{\rm GeV}^{-1}$.
We use Eq.~(\ref{alphaS-mine}) to evaluate spectrum of light hadrons listed in Table \ref{tab:lighthadron} with original results along side.
Although the errors seem to be comparatively greater while bag radius of proton is no longer $5\,{\rm GeV}^{-1}$,
a running $\alpha_{s}$ can remarkably reduce the errors of heavy hadrons, especially charmonium and bottomonium according to our calculation.

\begin{figure}[!ht]
    \centering
    \begin{tabular}{c}
        \includegraphics[width=0.4\textwidth]{Figures/AlphaS.eps}
    \end{tabular}
    \caption{Four running behaviors of strong coupling constant. 
    Bag radius $R$ takes range from $3\,{\rm GeV}^{-1}$ to $6\,{\rm GeV}^{-1}$, and the standard radius is $5\,{\rm GeV}^{-1}$ ($\approx 1\,{\rm fm}$) for checking.
    The solid line represents our result (\ref{alphaS-mine}). Dashed line draws Eq.~(\ref{alphaS}) when $\gamma=2.847$ and $\Lambda_{QCD}=0.281\,{\rm GeV}$.
    Dotted line draws Eq.~(\ref{alphaS}) when $\gamma=1$ and $\Lambda_{QCD}=0.281\,{\rm GeV}$. The dotdashed line indicates Eq.~(\ref{alphaS-2004}).}
    \label{fig:alphaS}
\end{figure}

\renewcommand{\tabcolsep}{0.39cm}
\renewcommand{\arraystretch}{1.3}
\begin{table}[!htbp]
    \caption{Comparations of bag radius $R$ (${\rm GeV}^{-1}$) and bag mass $M_{bag}$ (${\rm GeV}$) between this work and original result for light hadrons.
    The experimental $M_{exp}$ (${\rm GeV}$) takes average of isospin.
    In this work, we choose $Z_{0}=1.83$ with little change to reduce errors of nucleon.}
    \label{tab:lighthadron}
    \begin{tabular}{cccccc}
        \hline\hline
        \multirow{2}{*}{State} & \multicolumn{2}{c}{[\rm ref]} & \multicolumn{2}{c}{\rm this\ work} & \multirow{2}{*}{$M_{exp}$} \\ 
        \cline{2-3} \cline{4-5} 
        & $R$ & $M_{bag}$ & $R$ & $M_{bag}$ \\ \hline
        $N$               &5.00      &0.938      &5.22      &0.932      &0.939  \\ 
        $\Delta$          &5.48      &1.233      &5.33      &1.241      &1.232  \\ 
        $\Lambda$         &4.95      &1.105      &5.26      &1.096      &1.116  \\ 
        $\Sigma$          &4.95      &1.144      &5.22      &1.137      &1.193  \\ 
        $\Xi$             &4.91      &1.289      &5.27      &1.282      &1.318  \\ 
        $\Sigma^{\ast}$   &5.43      &1.382      &5.38      &1.383      &1.385  \\ 
        $\Xi^{\ast}$      &5.39      &1.529      &5.42      &1.529      &1.533  \\ 
        $\Omega$          &5.35      &1.672      &5.46      &1.677      &1.672  \\ 
        $\pi$             &3.34      &0.280      &4.31      &0.348      &0.137  \\ 
        $\omega$          &4.71      &0.783      &4.55      &0.776      &0.783  \\ 
        $K$               &3.26      &0.497      &4.34      &0.561      &0.496  \\ 
        $K^{\ast}$        &4.65      &0.928      &4.63      &0.918      &0.894  \\ 
        $\phi$            &4.61      &1.068      &4.70      &1.064      &1.019  \\ 
        \hline\hline
    \end{tabular}
\end{table}

We now determine the masses of c, b quarks and five kinds of binding energy.
In this work, doubly heavy mesons are used for obtaining binding energy, by which we are unable to fit $m_{c}$ and $m_{b}$.
An alternative approach chooses heavy-light mesons like $D$, $D^{\ast}$, $B$, $B^{\ast}$. 
Additionally, we find the meson states with $J^{P}=1^{-}$ being easier to lie on zero-error level, while singly heavy baryons show stable errors.
It can be checked out in Table \ref{tab:lightmeson}, \ref{tab:lightbaryon} and \ref{tab:1heavybaryon}.
We thus use $D^{\ast}$ and $B^{\ast}$ states to give
\begin{equation}\label{mcmb}
    \begin{Bmatrix}
        m_{c}=1.641\,{\rm GeV} & m_{b}=5.093\,{\rm GeV}
    \end{Bmatrix}
\end{equation}
Having masses of c, b quarks, we then determine binding energy by $D_{s}^{\ast}$, $B_{s}^{\ast}$, $J/\psi$, $\Upsilon$ and $B_{c}^{\ast}$ giving
\begin{equation}\label{bingdings}
    \begin{Bmatrix}
        B_{cs}=-0.025\,{\rm GeV} & B_{cc}=-0.077\,{\rm GeV} \\ 
        B_{bs}=-0.032\,{\rm GeV} & B_{bb}=-0.128\,{\rm GeV} \\ 
        B_{bc}=-0.101\,{\rm GeV}
    \end{Bmatrix}
\end{equation}
where the mass $6.332\,{\rm GeV}$ of $B_{c}^{\ast}$ quotes Ref.~[ref].

In this work, the binding energy $B_{QQ^{\prime}}$ takes value for color $\boldsymbol{\bar{3}}$, which scales by Eq.~(\ref{colorfc}) through color configurations.
For color $\boldsymbol{1}$, we have $2B_{QQ^{\prime}}$. For color $\boldsymbol{6}$, we have $-B_{QQ^{\prime}}/2$.
Moreover, binding energies between quark and antiquark within tetraquark are $5B_{QQ^{\prime}}/4$ and $B_{QQ^{\prime}}/2$ for $\boldsymbol{6}\otimes\boldsymbol{\bar{6}}$ and $\boldsymbol{\bar{3}}\otimes\boldsymbol{3}$ respectively, 
whose scale can be found in Eq.~(\ref{cfcT}).

We have listed all the confirmed, lowest-lying and non-mixing hadronic states with magnetic moments and charge radii, in Table \ref{tab:lightmeson}, \ref{tab:heavymeson}, \ref{tab:lightbaryon} and \ref{tab:1heavybaryon}.
We also calculate doubly and triply heavy baryons presented in Table \ref{tab:23heavybaryon}.
There are short summaries:

1. Some theoretical masses $M_{bag}$ depart from the experimental values about $30\,{\rm GeV}-40\,{\rm GeV}$. 
Although they are not small errors, we still accept them up to the level following MIT bag model, in Table \ref{tab:lighthadron}, saying around $40\,{\rm GeV}$.

2. Table \ref{tab:lightmeson} and \ref{tab:heavymeson} list mesons without their antiparticles because they share the same mass, charge radius but minus magnetic moment.
Antiparticles can be neglected in studying tetraquarks as well.

3. Generally, MIT bag model discusses the magnetic moment ratio of hadron to proton, which is regarded to be more precise than magnetic moment itself.
In Table \ref{tab:lightbaryon}, the predicted ratios are close to experimental data, for which we could further calculate heavy baryons and tetraquarks well.

4. Our results of masses for $\Xi_{cc}$ and $\Xi_{cc}^{\ast}$ in Table \ref{tab:23heavybaryon} are $3.604\,{\rm GeV}$ and $3.714\,{\rm GeV}$ respectively.
We note that the errors of excited states are usually little in strange and charmed baryons in this work. Considering that Elbert has given $M\left(\Xi_{cc}^{\ast}\right)=3.727\,{\rm GeV}$, we believe $3.714\,{\rm GeV}$ will be a close prediction.

5. We calculate the confirmed states not for restudying mesons and baryons, but to prove that MIT bag model can determine spectrum well, from which we can further study multi-heavy baryons and tetraquarks.

\renewcommand{\tabcolsep}{0.37cm}
\renewcommand{\arraystretch}{1.3}
\begin{table}[!htbp]
    \caption{Spectrum of lowest-lying light mesons. $\mu_{p}$ is theoretical magnetic moment of proton. Charge radius is in unit of fm.
    The blank cells follow above since MIT bag model confuses with isospin.}
    \label{tab:lightmeson}
    \begin{tabular}{cccccc}
        \hline\hline
        {\rm State}         &$R_{0}$    &$M_{bag}$  &$M_{exp}$  &$\mu_{bag}/\mu_{p}$    &$r_{E}$ \\ \hline
        $\pi^{+}$ 	        &4.31	    &0.348	    &0.137	    &-	                    &0.62 \\ 
        $\omega$ 	        &4.55	    &0.776	    &0.783	    &0                      &0 \\ 
        $K^{+}$ 	        &4.34	    &0.561	    &0.494	    &-                      &0.61 \\ 
        $K^{0}$ 		    &           &           &0.498	    &-	                    &0.13 \\ 
        $K^{\ast +}$ 	    &4.63	    &0.918	    &0.892	    &0.82                   &0.65 \\ 
        $K^{\ast 0}$ 	    &           &           &0.896	    &-0.06	                &0.15 \\ 
        $\phi$ 	            &4.70	    &1.064	    &1.019	    &0                      &0 \\ 
        \hline\hline
    \end{tabular}
\end{table}

\renewcommand{\tabcolsep}{0.3cm}
\renewcommand{\arraystretch}{1.3}
\begin{table}[!htbp]
    \caption{Spectrum of lowest-lying heavy mesons. The blank cell follows $5.325\,{\rm GeV}$ because it doesn’t identify specific state.}
    \label{tab:heavymeson}
    \begin{tabular}{cccccc}
        \hline\hline
        {\rm State}             &$R_{0}$    &$M_{bag}$      &$M_{exp}$  &$\mu_{bag}/\mu_{p}$    &$r_{E}$ \\ \hline
        $D^{+}$ 	            &3.63	    &1.835	        &1.870	    &-                      &0.46 \\ 
        $D^{0}$ 			    &           &               &1.865	    &-                      &0.25 \\ 
        $D^{\ast +}$ 	        &4.09	    &2.009(fitted)  &2.010	    &0.43                   &0.51 \\ 
        $D^{\ast 0}$ 		    &           &               &2.007	    &-0.35                  &0.29 \\ 
        $B^{+}$ 	            &3.14	    &5.248	        &5.279	    &-                      &0.42 \\ 
        $B^{0}$ 			    &           &               &5.280	    &-                      &0.17 \\ 
        $B^{\ast +}$ 	        &3.47	    &5.325(fitted)  &5.325      &0.47                   &0.46 \\ 
        $B^{\ast 0}$ 		    &           &               &           &-0.19                  &0.19 \\ 
        $D_{s}^{+}$ 	        &3.77	    &1.961	        &1.968	    &-                      &0.46 \\ 
        $D_{s}^{\ast +}$ 	    &4.17	    &2.112(fitted)  &2.112	    &0.39                   &0.51 \\ 
        $B_{s}^{0}$ 	        &3.35	    &5.346	        &5.367	    &-                      &0.16 \\ 
        $B_{s}^{\ast 0}$ 	    &3.62	    &5.415(fitted)  &5.415	    &0.36                   &0.17 \\ 
        $\eta_{c}$ 	            &3.15	    &3.002	        &2.984	    &-                      &0 \\ 
        $J/\psi$ 	            &3.54	    &3.097(fitted)  &3.097	    &0                      &0 \\ 
        $B_{c}^{+}$ 	        &2.53	    &6.274	        &6.275	    &-                      &0.29 \\ 
        $B_{c}^{\ast +}$ 	    &2.81	    &6.332(fitted)  &6.332	    &0.19                   &0.32 \\ 
        $\eta_{b}$ 	            &1.59	    &9.397	        &9.399	    &-                      &0 \\ 
        $\Upsilon$ 	            &1.80	    &9.460(fitted)  &9.460	    &0                      &0 \\ 
        \hline\hline
    \end{tabular}
\end{table}

\renewcommand{\tabcolsep}{0.25cm}
\renewcommand{\arraystretch}{1.3}
\begin{table}[!htbp]
    \caption{Spectrum of lowest-lying light baryons. $\mu_{p}$ dividing $\mu_{bag}$ is theoretical value and another one beside is experimental data.}
    \label{tab:lightbaryon}
    \begin{tabular}{ccccccc}
        \hline\hline
        {\rm State}         &$R_{0}$    &$M_{bag}$  &$M_{exp}$  &$\mu_{bag}/\mu_{p}$    &$\mu_{exp}/\mu_{p}$    &$r_{E}$ \\ \hline
        $p$ 	            &5.22	    &0.932	    &0.938	    &1	                    &1	                    &0.75 \\ 
        $n$ 			    &           &           &0.940	    &-0.67	                &-0.69	                &0 \\ 
        $\Delta^{++}$ 	    &5.33	    &1.241	    &1.231	    &2.04	                &2.20	                &1.08 \\ 
        $\Delta^{+}$ 	    &           &           &1.235	    &1.02	                &0.97	                &0.77 \\ 
        $\Delta^{0}$ 		&           &           &1.233	    &0	                    &-	                    &0 \\ 
        $\Lambda$ 	        &5.26	    &1.096	    &1.116	    &-0.26	                &-0.22	                &0.18 \\ 
        $\Sigma^{+}$ 	    &5.22	    &1.137	    &1.189	    &0.98	                &0.88	                &0.77 \\ 
        $\Sigma^{0}$ 		&           &           &1.193	    &0.31	                &-	                    &0.17 \\ 
        $\Sigma^{-}$ 		&           &           &1.197	    &-0.36	                &-0.42	                &0.73 \\ 
        $\Sigma^{\ast +}$ 	&5.38	    &1.383	    &1.383	    &1.11	                &-	                    &0.79 \\ 
        $\Sigma^{\ast 0}$   &           &           &1.384	    &0.08	                &-	                    &0.18 \\ 
        $\Sigma^{\ast -}$ 	&           &           &1.387	    &-0.95	                &-	                    &0.75 \\ 
        $\Xi^{0}$ 	        &5.27       &1.282	    &1.315	    &-0.57	                &-0.45	                &0.25 \\ 
        $\Xi^{-}$ 			&           &           &1.322	    &-0.23	                &-0.23	                &0.72 \\ 
        $\Xi^{\ast 0}$ 	    &5.42	    &1.529	    &1.532	    &0.17	                &-	                    &0.26 \\ 
        $\Xi^{\ast -}$ 		&           &           &1.535	    &-0.87	                &-	                    &0.74 \\ 
        $\Omega$ 	        &5.46	    &1.677	    &1.672	    &-0.79	                &-0.72	                &0.72 \\ 
        \hline\hline
    \end{tabular}
\end{table}

\renewcommand{\tabcolsep}{0.37cm}
\renewcommand{\arraystretch}{1.3}
\begin{table}[!htbp]
    \caption{Spectrum of lowest-lying and non-mixing singly heavy baryons.}
    \label{tab:1heavybaryon}
    \begin{tabular}{cccccc}
        \hline\hline
        {\rm State}             &$R_{0}$    &$M_{bag}$  &$M_{exp}$  &$\mu_{bag}/\mu_{p}$    &$r_{E}$ \\ \hline
        $\Lambda_{c}$ 	        &4.86	    &2.270	    &2.286	    &0.18	                &0.60 \\ 
        $\Sigma_{c}^{++}$ 	    &4.82	    &2.411	    &2.454	    &0.76	                &0.92 \\ 
        $\Sigma_{c}^{+}$ 		&           &           &2.453	    &0.15	                &0.60 \\ 
        $\Sigma_{c}^{0}$ 		&           &           &2.454	    &-0.47	                &0.35 \\ 
        $\Sigma_{c}^{\ast ++}$ 	&5.01	    &2.512	    &2.518	    &1.46	                &0.95 \\ 
        $\Sigma_{c}^{\ast +}$ 	&           &           &2.518	    &0.50	                &0.62 \\ 
        $\Sigma_{c}^{\ast 0}$ 	&           &           &2.518	    &-0.46	                &0.37 \\ 
        $\Omega_{c}$ 	        &4.93	    &2.680	    &2.695	    &-0.38                  &0.28 \\ 
        $\Omega_{c}^{\ast}$ 	&5.10	    &2.764	    &2.766	    &-0.32                  &0.29 \\ 
        $\Lambda_{b}$ 	        &4.60	    &5.648	    &5.620	    &-0.03	                &0.25 \\ 
        $\Sigma_{b}^{+}$ 	    &4.64	    &5.835	    &5.811	    &0.80	                &0.72 \\ 
        $\Sigma_{b}^{0}$ 		&           &           &-	        &0.21	                &0.26 \\ 
        $\Sigma_{b}^{-}$ 	    &           &           &5.816	    &-0.39	                &0.62 \\ 
        $\Sigma_{b}^{\ast +}$ 	&4.73	    &5.872	    &5.830	    &1.18	                &0.73 \\ 
        $\Sigma_{b}^{\ast 0}$ 	&           &           &-	        &0.27	                &0.26 \\ 
        $\Sigma_{b}^{\ast -}$ 	&           &           &5.835	    &-0.64	                &0.63 \\  
        $\Omega_{b}$ 	        &4.77	    &6.080	    &6.046	    &-0.31                  &0.60 \\ 
        $\Omega_{b}^{\ast}$ 	&4.84	    &6.112	    &-	        &-0.51                  &0.60 \\ 
        \hline\hline
    \end{tabular}
\end{table}

\renewcommand{\tabcolsep}{0.51cm}
\renewcommand{\arraystretch}{1.3}
\begin{table}[!htbp]
    \caption{Spectrum of lowest-lying and non-mixing doubly and triply heavy baryons.}
    \label{tab:23heavybaryon}
    \begin{tabular}{ccccc}
        \hline\hline
        {\rm State}             &$R_{0}$    &$M_{bag}$  &$\mu_{bag}/\mu_{p}$    &$r_{E}$ \\ \hline          
        $\Xi_{cc}^{++}$ 	    &4.42	    &3.604	    &0.05                   &0.78 \\ 
        $\Xi_{cc}^{+}$ 			&           &           &0.33	                &0.45 \\ 
        $\Xi_{cc}^{\ast ++}$ 	&4.64	    &3.714	    &0.94	                &0.82 \\ 
        $\Xi_{cc}^{\ast +}$ 	&           &           &0.06	                &0.47 \\ 
        $\Xi_{bb}^{0}$ 	        &3.71	    &10.311	    &-0.20	                &0.29 \\ 
        $\Xi_{bb}^{-}$ 			&           &           &0.04	                &0.45 \\ 
        $\Xi_{bb}^{\ast 0}$ 	&3.87	    &10.360	    &0.43	                &0.30 \\ 
        $\Xi_{bb}^{\ast -}$ 	&           &           &-0.31	                &0.47 \\ 
        $\Omega_{cc}$ 	        &4.49	    &3.726	    &0.31	                &0.48 \\ 
        $\Omega_{cc}^{\ast}$ 	&4.69	    &3.820	    &0.12	                &0.50 \\ 
        $\Omega_{bb}$ 	        &3.83	    &10.408	    &0.03	                &0.45 \\ 
        $\Omega_{bb}^{\ast}$ 	&3.97	    &10.451	    &-0.27	                &0.47 \\ 
        $\Omega_{ccc}$ 	        &4.22	    &4.841	    &0.52	                &0.68 \\ 
        $\Omega_{ccb}$ 	        &3.75	    &8.112	    &0.24	                &0.44 \\ 
        $\Omega_{ccb}^{\ast}$ 	&3.83	    &8.133	    &0.31	                &0.45 \\ 
        $\Omega_{cbb}$ 	        &3.18	    &11.373	    &-0.10	                &0.11 \\ 
        $\Omega_{cbb}^{\ast}$ 	&3.31	    &11.402	    &0.10	                &0.11 \\ 
        $\Omega_{bbb}$ 	        &2.59	    &14.626	    &-0.09	                &0.28 \\ 
        \hline\hline
    \end{tabular}
\end{table}


\section{Calculations of Mixing States and Tetraquarks}\label{Hadrons}
\subsection{Baryons of Mixing Systems}
\renewcommand{\tabcolsep}{0.73cm}
\renewcommand{\arraystretch}{1.3}
\begin{table*}[!htbp]
    \caption{Spectrum of lowest-lying and mixing heavy baryons. The experimental data $M_{exp}$ takes average of isospin.
    Magnetic moments and charge radii are organized in the order of $I_{3}=\frac{1}{2},\,-\frac{1}{2}$ for $I=\frac{1}{2}$.}
    \label{tab:mixingbaryon}
    \begin{tabular}{ccccccc}
        \hline\hline
        {\rm State}         &{\rm Eigenvector}  &$R_{0}$    &$M_{bag}$  &$M_{exp}$  &$\mu_{bag}/\mu_{p}$    &$r_{E}$ \\ \hline
        \multirow{2}{*}{$\left(\Xi_{c}^{\prime},\ \Xi_{c}\right)$}
                            &(0.05, 1.00)       &4.89	    &2.436	    &2.470      &0.18, 0.18             &0.63, 0.32 \\ 
                            &(-1.00, 0.05)      &4.88	    &2.544	    &2.579      &0.20, -0.43            &0.63, 0.32 \\ 
        $\Xi_{c}^{\ast}$ 	&1.00               &5.06	    &2.636	    &2.646      &0.58, -0.39            &0.65, 0.33 \\ 
        \multirow{2}{*}{$\left(\Xi_{b}^{\prime},\ \Xi_{b}\right)$}
                            &(0.01, 1.00)       &4.64	    &5.805	    &5.795      &-0.03, -0.03           &0.30, 0.60 \\ 
                            &(-1.00, 0.01)      &4.71	    &5.956	    &5.935      &0.25, -0.35            &0.30, 0.61 \\ 
        $\Xi_{b}^{\ast}$ 	&1.00               &4.79	    &5.991	    &5.954      &0.34, -0.58            &0.31, 0.62 \\ 
        \multirow{2}{*}{$\left(\Xi_{bc}^{\prime},\ \Xi_{bc}\right)$}
                            &(0.39, 0.92)       &4.22	    &7.015	    &-          &0.45, -0.20            &0.58, 0.19 \\ 
                            &(-0.92, 0.39)      &4.09	    &6.953	    &-          &0.01, 0.12             &0.56, 0.19 \\ 
        $\Xi_{bc}^{\ast}$ 	&1.000              &4.31	    &7.044	    &-          &0.69, -0.13            &0.59, 0.20 \\ 
        \multirow{2}{*}{$\left(\Omega_{bc}^{\prime},\ \Omega_{bc}\right)$}
                            &(0.40, 0.92)       &4.29	    &7.116	    &-          &-0.16                  &0.15 \\ 
                            &(-0.92, 0.40)      &4.18	    &7.064	    &-          &0.11                   &0.14 \\ 
        $\Omega_{bc}^{\ast}$&1.000              &4.37	    &7.142	    &-          &-0.08                  &0.15 \\ 
        \hline\hline
    \end{tabular}
\end{table*}

\renewcommand{\tabcolsep}{0.62cm}
\renewcommand{\arraystretch}{1.3}
\begin{table*}[!htbp]
    \caption{Spectrum of singly heavy tetraquarks $nn\bar{s}\bar{c}$ and $nn\bar{s}\bar{b}$. 
    Magnetic moments and charge radii are organized in the order of $I_{3}=1,\,0,\,-1$ for $I=1$.}
    \label{tab:1heavytetraquark}
    \begin{tabular}{ccccccc}
        \hline\hline
        {\rm State}     &$J^{P}$    &{\rm Eigenvector}      &$R_{0}$    &$M_{bag}$  &$\mu_{bag}/\mu_{p}$    &$r_{E}$ \\ \hline
        \multirow{2}{*}{${\left(nn\bar{s}\bar{c}\right)}^{I=1}$}
                        &\multirow{2}{*}{$0^{+}$}
                                    &(0.54, 0.84)           &5.73	    &3.218	    &-                      &0.91, 0.38, 0.73 \\ 
                        &           &(-0.84, 0.55)          &5.39	    &2.776	    &-                      &0.85, 0.35, 0.69 \\ 
                        &\multirow{3}{*}{$1^{+}$}
                                    &(0.81, 0.58, 0.10)     &5.46	    &3.001	    &0.95, 0.26, -0.43      &0.86, 0.36, 0.70 \\ 
                        &           &(0.25, -0.49, 0.84)    &5.57	    &3.154	    &0.44, 0.16, -0.13      &0.88, 0.37, 0.71 \\ 
                        &           &(-0.54, 0.65, 0.54)    &5.38	    &2.846	    &0.60, 0.43, -0.57      &0.85, 0.35, 0.69 \\ 
                        &$2^{+}$    &1.00                   &5.64	    &3.075	    &1.53, 0.45, -0.63      &0.89, 0.37, 0.72 \\ 
        \multirow{2}{*}{${\left(nn\bar{s}\bar{c}\right)}^{I=0}$}
                        &\multirow{2}{*}{$0^{+}$}
                                    &(0.63, 0.77)           &5.56	    &2.934	    &-                      &0.37 \\ 
                        &           &(-0.78, 0.63)          &5.19	    &2.513	    &-                      &0.34 \\ 
                        &\multirow{3}{*}{$1^{+}$}
                                    &(0.77, 0.07, 0.64)     &5.40	    &2.895	    &0.16                   &0.35 \\ 
                        &           &(-0.23, -0.90, 0.36)   &5.57	    &3.056	    &0.31                   &0.37 \\ 
                        &           &(-0.60, 0.43, 0.68)    &5.35	    &2.674	    &0.17                   &0.35 \\ 
                        &$2^{+}$    &1.00                   &5.66	    &3.063	    &0.45                   &0.37 \\ 

        \multirow{2}{*}{${\left(nn\bar{s}\bar{b}\right)}^{I=1}$}
                        &\multirow{2}{*}{$0^{+}$}
                                    &(0.53, 0.85)           &5.53	    &6.580	    &-                      &1.07, 0.71, 0.36 \\ 
                        &           &(-0.85, 0.53)          &5.28	    &6.202	    &-                      &1.02, 0.68, 0.34 \\ 
                        &\multirow{3}{*}{$1^{+}$}
                                    &(0.66, 0.75, 0.06)     &5.35	    &6.419	    &1.13, 0.33, -0.47      &1.03, 0.69, 0.35 \\ 
                        &           &(0.36, -0.38, 0.85)    &5.43	    &6.554	    &0.52, 0.30, 0.09       &1.05, 0.70, 0.35 \\ 
                        &           &(-0.66, 0.55, 0.52)    &5.22	    &6.228	    &0.83, 0.31, -0.21      &1.01, 0.67, 0.34 \\ 
                        &$2^{+}$    &1.00                   &5.47	    &6.446	    &1.69, 0.64, -0.40      &1.05, 0.70, 0.36 \\ 
        \multirow{2}{*}{${\left(nn\bar{s}\bar{b}\right)}^{I=0}$}
                        &\multirow{2}{*}{$0^{+}$}
                                    &(0.66, 0.75)           &5.41	    &6.327	    &-                      &0.69 \\ 
                        &           &(-0.75, 0.66)          &5.14	    &5.980	    &-                      &0.66 \\ 
                        &\multirow{3}{*}{$1^{+}$}                            
                                    &(0.66, -0.23, 0.72)    &5.33	    &6.322	    &0.30                   &0.68 \\ 
                        &           &(-0.46, -0.88, 0.14)   &5.38	    &6.454	    &0.34                   &0.69 \\ 
                        &           &(-0.60, 0.43, 0.68)    &5.14	    &6.038	    &0.30                   &0.66 \\ 
                        &$2^{+}$    &1.00                   &5.48	    &6.431	    &0.64                   &0.70 \\ 
        \hline\hline
    \end{tabular}
\end{table*}

\renewcommand{\tabcolsep}{0.59cm}
\renewcommand{\arraystretch}{1.2}
\begin{table*}[!htbp]
    \caption{Spectrum of doubly heavy tetraquarks $nn\bar{c}\bar{c}$, $nn\bar{b}\bar{b}$ and $nn\bar{c}\bar{b}$. 
    Magnetic moments and charge radii are organized in the order of $I_{3}=1,\,0,\,-1$ for $I=1$.}
    \label{tab:2heavytetraquark1}
    \begin{tabular}{ccccccc}
        \hline\hline
        {\rm State}     &$J^{P}$    &{\rm Eigenvector}      &$R_{0}$    &$M_{bag}$  &$\mu_{bag}/\mu_{p}$    &$r_{E}$ \\ \hline
        \multirow{2}{*}{${\left(nn\bar{c}\bar{c}\right)}^{I=1}$}
                        &\multirow{2}{*}{$0^{+}$}
                                    &(0.40, 0.92)           &5.40	    &4.342	    &-                      &0.56, 0.54, 0.94 \\ 
                        &           &(-0.91, 0.41)          &5.04	    &4.032	    &-                      &0.52, 0.50, 0.88 \\ 
                        &$1^{+}$    &1.00                   &5.22	    &4.117	    &0.49, -0.02, -0.51     &0.54, 0.52, 0.91 \\ 
                        &$2^{+}$    &1.00                   &5.32	    &4.179	    &1.00, -0.02, -1.04     &0.55, 0.53, 0.93 \\ 
        \multirow{2}{*}{${\left(nn\bar{c}\bar{c}\right)}^{I=0}$}
                        &\multirow{2}{*}{$1^{+}$}
                                    &(0.97, 0.25)           &5.15	    &3.925	    &-0.32                  &0.51 \\ 
                        &           &(-0.24, 0.97)          &5.30	    &4.205	    &0.30                   &0.53 \\ 
        \multirow{2}{*}{${\left(nn\bar{b}\bar{b}\right)}^{I=1}$}
                        &\multirow{2}{*}{$0^{+}$}
                                    &(0.17, 0.99)           &4.90	    &11.093	    &-                      &0.92, 0.59, 0.38 \\ 
                        &           &(-0.98, 0.18)          &4.77	    &10.834	    &-                      &0.90, 0.58, 0.37 \\ 
                        &$1^{+}$    &1.00                   &4.83	    &10.854	    &0.65, 0.18, -0.28      &0.91, 0.58, 0.38 \\ 
                        &$2^{+}$    &1.00                   &4.88	    &10.878	    &1.31, 0.37, -0.56      &0.92, 0.59, 0.38 \\ 
        \multirow{2}{*}{${\left(nn\bar{b}\bar{b}\right)}^{I=0}$}
                        &\multirow{2}{*}{$1^{+}$}
                                    &(1.00, 0.08)           &4.76	    &10.654	    &0.06                   &0.57 \\ 
                        &           &(-0.08, 1.00)          &4.83	    &10.982	    &0.31                   &0.58 \\ 

        \multirow{2}{*}{${\left(nn\bar{c}\bar{b}\right)}^{I=1}$}
                        &\multirow{2}{*}{$0^{+}$}
                                    &(0.30, 0.95)           &5.16	    &7.714	    &-                      &0.78, 0.25, 0.70 \\ 
                        &           &(-0.95, 0.31)          &4.91	    &7.438	    &-                      &0.74, 0.23, 0.67 \\ 
                        &\multirow{3}{*}{$1^{+}$}
                                    &(0.65, 0.76, 0.09)     &5.04	    &7.509	    &0.98, 0.22, -0.54      &0.76, 0.24, 0.68 \\ 
                        &           &(0.13, -0.22, 0.97)    &5.10	    &7.699	    &-0.07, -0.12, -0.18    &0.77, 0.24, 0.69 \\ 
                        &           &(-0.75, 0.61, 0.24)    &4.96	    &7.465	    &0.79, 0.16, -0.47      &0.75, 0.24, 0.67 \\ 
                        &$2^{+}$    &1.00                   &5.12	    &7.531	    &1.16, 0.18, -0.80      &0.78, 0.25, 0.69 \\ 
        \multirow{2}{*}{${\left(nn\bar{c}\bar{b}\right)}^{I=0}$}
                        &\multirow{2}{*}{$0^{+}$}
                                    &(0.93, 0.37)           &4.96	    &7.502	    &-                      &0.24 \\ 
                        &           &(-0.38, 0.93)          &4.84	    &7.260	    &-                      &0.23 \\ 
                        &\multirow{3}{*}{$1^{+}$}
                                    &(0.91, -0.36, 0.21)    &4.93	    &7.518	    &0.10                   &0.23 \\ 
                        &           &(-0.39, -0.92, 0.10)   &5.07	    &7.605	    &0.28                   &0.24 \\ 
                        &           &(-0.16, 0.18, 0.97)    &4.90	    &7.288	    &-0.13                  &0.23 \\ 
                        &$2^{+}$    &1.00                   &5.14	    &7.483	    &0.18                   &0.25 \\ 
        \hline\hline
    \end{tabular}
\end{table*}

\renewcommand{\tabcolsep}{0.84cm}
\renewcommand{\arraystretch}{1.2}
\begin{table*}[!htbp]
    \caption{Spectrum of doubly heavy tetraquarks $ss\bar{c}\bar{c}$, $ss\bar{b}\bar{b}$ and $ss\bar{c}\bar{b}$.}
    \label{tab:2heavytetraquark0}
    \begin{tabular}{ccccccc}
        \hline\hline
        {\rm State}     &$J^{P}$    &{\rm Eigenvector}      &$R_{0}$    &$M_{bag}$  &$\mu_{bag}/\mu_{p}$    &$r_{E}$ \\ \hline
        \multirow{2}{*}{$ss\bar{c}\bar{c}$}
                        &\multirow{2}{*}{$0^{+}$}
                                    &(0.49, 0.87)           &5.48	    &4.521	    &-                      &0.92 \\ 
                        &           &(-0.87, 0.50)          &5.13	    &4.299	    &-                      &0.87 \\ 
                        &$1^{+}$    &1.00                   &5.30	    &4.382	    &-0.44                  &0.89 \\ 
                        &$2^{+}$    &1.00                   &5.39	    &4.433	    &-0.88                  &0.91 \\ 
        \multirow{2}{*}{$ss\bar{b}\bar{b}$}
                        &\multirow{2}{*}{$0^{+}$}
                                    &(0.24, 0.97)           &5.01	    &11.232	    &-                      &0.32 \\ 
                        &           &(-0.97, 0.25)          &4.88	    &11.078	    &-                      &0.31 \\ 
                        &$1^{+}$    &1.00                   &4.94	    &11.098	    &-0.21                  &0.31 \\ 
                        &$2^{+}$    &1.00                   &4.98	    &11.119	    &-0.43                  &0.32 \\ 
        \multirow{2}{*}{$ss\bar{c}\bar{b}$}
                        &\multirow{2}{*}{$0^{+}$}
                                    &(0.40, 0.92)           &5.26	    &7.875	    &-                      &0.67 \\ 
                        &           &(-0.91, 0.40)          &5.01	    &7.693	    &-                      &0.64 \\ 
                        &\multirow{3}{*}{$1^{+}$}
                                    &(0.70, 0.71, 0.11)     &4.98	    &7.757	    &-0.41                  &0.64 \\ 
                        &           &(0.17, -0.32, 0.93)    &5.00	    &7.858	    &-0.17                  &0.64 \\ 
                        &           &(-0.69, 0.63, 0.35)    &4.88	    &7.716	    &-0.24                  &0.62 \\ 
                        &$2^{+}$    &1.00                   &5.20	    &7.779	    &-0.66                  &0.66 \\ 
        \hline\hline
    \end{tabular}
\end{table*}

\renewcommand{\tabcolsep}{0.74cm}
\renewcommand{\arraystretch}{1.3}
\begin{table*}[!htbp]
    \caption{Spectrum of doubly heavy tetraquarks $ns\bar{c}\bar{c}$, $ns\bar{b}\bar{b}$. 
    Magnetic moments and charge radii are organized in the order of $I_{3}=\frac{1}{2},\,-\frac{1}{2}$ for $I=\frac{1}{2}$.}
    \label{tab:2heavytetraquark1/2}
    \begin{tabular}{ccccccc}
        \hline\hline
        {\rm State}     &$J^{P}$    &{\rm Eigenvector}      &$R_{0}$    &$M_{bag}$  &$\mu_{bag}/\mu_{p}$    &$r_{E}$ \\ \hline
        \multirow{2}{*}{$ns\bar{c}\bar{c}$}
                        &\multirow{2}{*}{$0^{+}$}
                                    &(0.44, 0.90)           &5.44	    &4.429	    &-                      &0.51, 0.93 \\ 
                        &           &(-0.89, 0.45)          &5.09	    &4.165	    &-                      &0.48, 0.88 \\ 
                        &\multirow{3}{*}{$1^{+}$}
                                    &(0.99, -0.07, 0.09)    &5.16	    &4.247	    &0.03, -0.47            &0.49, 0.89 \\ 
                        &           &(-0.11, -0.28, 0.95)   &5.23	    &4.314	    &0.35, -0.57            &0.49, 0.90 \\ 
                        &           &(0.04, 0.96, 0.29)     &5.04	    &4.091	    &-0.29, -0.37           &0.48, 0.87 \\ 
                        &$2^{+}$    &1.00                   &5.36	    &4.305	    &0.07, -0.96            &0.50, 0.92 \\ 
        \multirow{2}{*}{$ns\bar{b}\bar{b}$}             
                        &\multirow{2}{*}{$0^{+}$}             
                                    &(0.20, 0.98)           &4.96	    &11.160	    &-                      &0.62, 0.35 \\ 
                        &           &(-0.98, 0.20)          &4.83	    &10.955	    &-                      &0.60, 0.34 \\ 
                        &\multirow{3}{*}{$1^{+}$}
                                    &(1.00, -0.01, 0.03)    &4.79	    &10.974	    &0.22, -0.24            &0.60, 0.34 \\ 
                        &           &(-0.03, -0.09, 1.00)   &4.77	    &11.068	    &0.37, -0.54            &0.60, 0.34 \\ 
                        &           &(0.01, 1.00, 0.10)     &4.66	    &10.811	    &0.06, 0.06             &0.58, 0.33 \\ 
                        &$2^{+}$    &1.00                   &4.93	    &10.997	    &0.45, -0.50            &0.62, 0.35 \\ 
        \hline\hline
    \end{tabular}
\end{table*}

Recall chromomagnetic mixing, that a diquark or antiquark carring different flavors beyond SU(2) symmetry, can be symmetric or asymmetric in terms of flavor wave function.
At this point, we would study the mixing systems $\left(\phi^{B}\chi_{2}^{B},\phi^{B}\chi_{3}^{B}\right)$ for baryons with $J^{P}={\frac{1}{2}}^{+}$, 
where $\Xi_{c}$ , $\Xi_{b}$, $\Xi_{bc}$, $\Omega_{bc}$ are involved.

We first present the matrix of CMI for this mixing system in Appendix C, where the subscripts of $C_{ij}$ depend on flavor configurations. 
In this paper, we propose $(ns)c$ for $\Xi_{c}$, $(ns)b$ for $\Xi_{b}$, $(bc)n$ for $\Xi_{bc}$ and $(bc)s$ for $\Omega_{bc}$, since two light or heavy quarks are regarded to be easier to form diquark.
We also consider binding energies $B_{cs}$, $B_{bs}$, $B_{bc}$ and will add them to mass formula directly.

As noted in Sect.~\ref{Method}, variational mothed requires the mass formula with analytical eigenvalues of CMI matrix.
One will not put a set of initially guessed $x_{i}$, $R$ into matrix (\ref{mixing23}) and perform numerical variation, 
because even a mixing hadronic state has its unique dynamical information.
Following this principle, we diagonalize the matrix (\ref{mixing23}) getting two eigenvalues to construct mass formula.
Once $x_{i}$, $R$ are determined, eigenvectors can be evaluated for identfying and describing mixing states.

Table \ref{tab:mixingbaryon} reveals our results of mixing baryons.
In the second column, we notice that the mixing effects are not significant, even neglectable for singly heavy baryons.
Therefore the masses of mixing states can be compared with experimental data of $\Xi_{c}$, $\Xi_{c}^{\prime}$, $\Xi_{b}$ and $\Xi_{b}^{\prime}$ as listed in the fifth column, whose errors are acceptable in MIT bag model.
On the other hand, mixing effects will be reduced as charm quarks are replaced by bottom quarks in $(ns)c$, or strange quarks are replaced by n quarks in $(bc)s$, according to our results.
We believe this is due to breaking of Heavy Quark Symmetry and small contribution from off-diagonal elements of matrix (\ref{mixing23}).


\subsection{Singly Heavy Tetraquarks}
We now move on to singly heavy tetraquarks. In 2020, LHCb has reported the first exotic state $X_{0}(2900)$ with open heavy flavor.
Karliner's team has interpreted it as an isosinglet tetraquark $cs\bar{u}\bar{d}$ in Ref.[]. We will also calculate $nn\bar{s}\bar{c}$ and $nn\bar{s}\bar{b}$ systems via MIT bag model, which occurs controllable errors for mesons and baryons.
The states with three n quarks shall be neglected for stability.

To study chromomagnetic mixing of $nn\bar{s}\bar{c}$ and $nn\bar{s}\bar{b}$, we introduce the mixing systems such as
$\left(\phi_{2}^{T}\chi_{3}^{T},\phi_{1}^{T}\chi_{6}^{T}\right)$ with $J^{P}=0^{+}$ and $I=1$,
$\left(\phi_{1}^{T}\chi_{3}^{T},\phi_{2}^{T}\chi_{6}^{T}\right)$ with $J^{P}=0^{+}$ and $I=0$,
$\left(\phi_{2}^{T}\chi_{2}^{T},\phi_{2}^{T}\chi_{4}^{T},\phi_{1}^{T}\chi_{5}^{T}\right)$ with $J^{P}=1^{+}$ and $I=1$,
$\left(\phi_{1}^{T}\chi_{2}^{T},\phi_{1}^{T}\chi_{4}^{T},\phi_{2}^{T}\chi_{5}^{T}\right)$ with $J^{P}=1^{+}$ and $I=0$.
Additionally, binding energies become diagonal matrix in mass formula, since the mixing systems contain two color configurations and spin states are orthorhombic.
One should diagonalize the sum of interaction matrices in calculating hadron mass.

Following variational principle, we will diagonalize second order matrix getting two analytical eigenvalues to construct mass formula as usual.
On the contrary, however, we are unable to perform the same procedure on three order matrix, whose eigenvalues are roots of a cubic equation.
At this point, we would rather scan three sets of $x_{i}$, $R$ for solving equation numerically to obtain minimal masses with three eigenvalues respectively.

Our results are shown in Table \ref{tab:1heavytetraquark}.


\subsection{Doubly Heavy Tetraquarks}
At the end of this work, we study doubly heavy tetraquarks with strange number $S=0$, $S=1$ and $S=2$.
There are mixing systems such as $\left(\phi_{2}^{T}\chi_{3}^{T},\phi_{1}^{T}\chi_{6}^{T}\right)$ with $J^{P}=0^{+}$ for isotriplet $nn\bar{Q}\bar{Q}$, $ss\bar{Q}\bar{Q}$ and $ns\bar{Q}\bar{Q}$,
$\left(\phi_{2}^{T}\chi_{5}^{T},\phi_{1}^{T}\chi_{4}^{T}\right)$ with $J^{P}=1^{+}$ for isosinglet $nn\bar{Q}\bar{Q}$,
and $\left(\phi_{2}^{T}\chi_{2}^{T},\phi_{2}^{T}\chi_{5}^{T},\phi_{1}^{T}\chi_{4}^{T}\right)$ with $J^{P}=1^{+}$ for $ns\bar{Q}\bar{Q}$.
$nn\bar{c}\bar{b}$ follows the same case of $nn\bar{s}\bar{c}$. We will neglect $ns\bar{c}\bar{b}$ for there is mixing among six spin-color states, which we cannot calculate even by numerical variation.


\section{Conclusions and Discussions}\label{Conclusions}
In this work, we have proposed ..., with many quantitative corrections remained to be explored.


\medskip
\textbf{ACKNOWLEDGMENTS}

D. J thanks Xiang Liu and A. Hosaka for discussions. D. J. is supported by
the National Natural Science Foundation of China under the no. 11565023 and
the Feitian Distinguished Professor Program of Gansu (2014-2016). C.P is
supported by the National Natural Science Foundation of China under Grants
No. 11965016, and by the projects funded by Science and Technology
Department of Qinghai Province (No. 2018-ZJ-971Q and No. 2019-ZJ-A10), the
Key Laboratory of IoT of Qinghai under grant (No. 2020-ZJ-Y16).

\medskip
\section*{Appendix A}\label{AppendixA}
\setcounter{equation}{0}
\renewcommand{\theequation}{A\arabic{equation}}
For meson, baryon, and tetraquark systems labeled $(q_{1}\bar{q}_{2})$ $(q_{1}q_{2})q_{3}$ and $(q_{1}q_{2})(\bar{q}_{3}\bar{q}_{4})$ respectively, 
the total color wave functions studied from SU(3) symmetry read
\begin{equation}\label{pcolorM}
    \phi^{M}=\frac{1}{\sqrt{3}} \left(r\bar{r}+g\bar{g}+b\bar{b} \right)
\end{equation}
\begin{equation}\label{pcolorB}
    \phi^{B}=\frac{1}{\sqrt{6}} \left(gbr-bgr+brg-rbg+rgb-grb \right)
\end{equation}
\begin{equation}\label{pcolorT}
    \begin{aligned}
        \phi_{1}^{T}&=\frac{1}{\sqrt{6}} \left(rr\bar{r}\bar{r}+gg\bar{g}\bar{g}+bb\bar{b}\bar{b} \right)
        +\frac{1}{2\sqrt{6}} \left(rb\bar{b}\bar{r}+br\bar{b}\bar{r} \right.\\ 
        &\left. +gr\bar{g}\bar{r}+rg\bar{g}\bar{r}+gb\bar{b}\bar{g}+bg\bar{b}\bar{g}+gr\bar{r}\bar{g}+rg\bar{r}\bar{g}+gb\bar{g}\bar{b} \right. \\
        &\left. +bg\bar{g}\bar{b}+rb\bar{r}\bar{b}+br\bar{r}\bar{b} \right) \\ 
        \phi_{2}^{T}&=\frac{1}{2\sqrt{3}} \left(rb\bar{b}\bar{r}-br\bar{b}\bar{r}-gr\bar{g}\bar{r}+rg\bar{g}\bar{r}+gb\bar{b}\bar{g}-bg\bar{b}\bar{g} \right.\\
        &\left. +gr\bar{r}\bar{g}-rg\bar{r}\bar{g}-gb\bar{g}\bar{b}+bg\bar{g}\bar{b}-rb\bar{r}\bar{b}+br\bar{r}\bar{b}\right)
    \end{aligned}
\end{equation}
associated with states (\ref{colorM}), (\ref{colorB}) and (\ref{colorT}).

The matrices of color factors are available by applying Eq.~(\ref{colorfc}) to wave functions, which give
\begin{equation}\label{cfcM}
    \left\langle \boldsymbol{\lambda_{1}}\cdot \boldsymbol{\lambda_{2}} \right\rangle=-\frac{16}{3}
\end{equation}
for meson in $\left(\phi^{M}\right)$,
\begin{equation}\label{cfcB}
    \left\langle \boldsymbol{\lambda_{1}}\cdot \boldsymbol{\lambda_{2}} \right\rangle=
    \left\langle \boldsymbol{\lambda_{1}}\cdot \boldsymbol{\lambda_{3}} \right\rangle=
    \left\langle \boldsymbol{\lambda_{2}}\cdot \boldsymbol{\lambda_{3}} \right\rangle=-\frac{8}{3}
\end{equation}
for baryon in $\left(\phi^{B}\right)$, and
\begin{equation}\label{cfcT}
    \begin{aligned}
        &\left\langle \boldsymbol{\lambda_{1}}\cdot \boldsymbol{\lambda_{2}} \right\rangle=
        \left\langle \boldsymbol{\lambda_{3}}\cdot \boldsymbol{\lambda_{4}} \right\rangle=
        \begin{bmatrix}
            \frac{4}{3} & 0 \\
            0 & -\frac{8}{3}
        \end{bmatrix} \\ 
        &\left\langle \boldsymbol{\lambda_{1}}\cdot \boldsymbol{\lambda_{3}} \right\rangle=
        \left\langle \boldsymbol{\lambda_{2}}\cdot \boldsymbol{\lambda_{4}} \right\rangle=
        \begin{bmatrix}
            -\frac{10}{3} & 2\sqrt{2} \\
            2\sqrt{2} & -\frac{4}{3}
        \end{bmatrix} \\ 
        &\left\langle \boldsymbol{\lambda_{1}}\cdot \boldsymbol{\lambda_{4}} \right\rangle=
        \left\langle \boldsymbol{\lambda_{2}}\cdot \boldsymbol{\lambda_{3}} \right\rangle=
        \begin{bmatrix}
            -\frac{10}{3} & -2\sqrt{2} \\
            -2\sqrt{2} & -\frac{4}{3}
        \end{bmatrix}
    \end{aligned}
\end{equation}
for tetraquark in $(\phi_{1}^{T},\phi_{2}^{T})$ space.


\medskip
\section*{Appendix B}\label{AppendixB}
\setcounter{equation}{0}
\renewcommand{\theequation}{B\arabic{equation}}
For meson, baryon, and tetraquark systems labeled $(q_{1}\bar{q}_{2})$ $(q_{1}q_{2})q_{3}$ and $(q_{1}q_{2})(\bar{q}_{3}\bar{q}_{4})$ respectively, 
the total spin wave functions collected by checking C-G coefficients read
\begin{equation}\label{pspinM}
    \chi_{1}^{M}=\uparrow\uparrow \quad \chi_{2}^{M}=\frac{1}{\sqrt{2}} \left(\uparrow\downarrow-\downarrow\uparrow\right)
\end{equation}
\begin{equation}\label{pspinB}
    \begin{aligned}
        \chi_{1}^{B}&=\uparrow\uparrow\uparrow \\ 
        \chi_{2}^{B}&=\sqrt{\frac{2}{3}}\uparrow\uparrow\downarrow-\frac{1}{\sqrt{6}} \left(\uparrow\downarrow\uparrow+\downarrow\uparrow\uparrow\right) \\ 
        \chi_{3}^{B}&=\frac{1}{\sqrt{2}} \left(\uparrow\downarrow\uparrow-\downarrow\uparrow\uparrow\right)
    \end{aligned}
\end{equation}
\begin{equation}\label{pspinT}
    \begin{aligned}
        \chi_{1}^{T}&=\uparrow\uparrow\uparrow\uparrow \\ 
        \chi_{2}^{T}&=\frac{1}{2} \left(\uparrow\uparrow\uparrow\downarrow+\uparrow\uparrow\downarrow\uparrow-\uparrow\downarrow\uparrow\uparrow-\downarrow\uparrow\uparrow\uparrow\right) \\ 
        \chi_{3}^{T}&=\frac{1}{\sqrt{3}} \left(\uparrow\uparrow\downarrow\downarrow+\downarrow\downarrow\uparrow\uparrow\right)\\
        &-\frac{1}{2\sqrt{3}} \left(\uparrow\downarrow\uparrow\downarrow+\uparrow\downarrow\downarrow\uparrow+\downarrow\uparrow\uparrow\downarrow+\downarrow\uparrow\downarrow\uparrow\right) \\ 
        \chi_{4}^{T}&=\frac{1}{\sqrt{2}} \left(\uparrow\uparrow\uparrow\downarrow-\uparrow\uparrow\downarrow\uparrow\right) \\ 
        \chi_{5}^{T}&=\frac{1}{\sqrt{2}} \left(\uparrow\downarrow\uparrow\uparrow-\downarrow\uparrow\uparrow\uparrow\right) \\ 
        \chi_{6}^{T}&=\frac{1}{2} \left(\uparrow\downarrow\uparrow\downarrow-\uparrow\downarrow\downarrow\uparrow-\downarrow\uparrow\uparrow\downarrow+\downarrow\uparrow\downarrow\uparrow\right)
    \end{aligned}
\end{equation}
associated with states (\ref{spinM}), (\ref{spinB}) and (\ref{spinT}).

The matrices of spin factors are available by applying Eq.~(\ref{spinfc}) to wave functions, which give
\begin{equation}\label{sfcM}
    \left\langle \boldsymbol{\sigma_{1}}\cdot \boldsymbol{\sigma_{2}} \right\rangle=
    \begin{bmatrix}
        1 & 0 \\
        0 & -3
    \end{bmatrix}
\end{equation}
for meson in $(\chi_{1}^{M},\chi_{2}^{M})$ space,
    \begin{align}
        \left\langle \boldsymbol{\sigma_{1}}\cdot \boldsymbol{\sigma_{2}} \right\rangle&=
        \begin{bmatrix}
            1 & 0 & 0 \\
            0 & 1 & 0 \\
            0 & 0 & -3
        \end{bmatrix} \\ 
        \left\langle \boldsymbol{\sigma_{1}}\cdot \boldsymbol{\sigma_{3}} \right\rangle&=
        \begin{bmatrix}
            1 & 0 & 0 \\
            0 & -2 & -\sqrt{3} \\
            0 & -\sqrt{3} & 0
        \end{bmatrix} \\ 
        \left\langle \boldsymbol{\sigma_{2}}\cdot \boldsymbol{\sigma_{3}} \right\rangle&=
        \begin{bmatrix}
            1 & 0 & 0 \\
            0 & -2 & \sqrt{3} \\
            0 & \sqrt{3} & 0
        \end{bmatrix}
    \end{align}
for baryon in $(\chi_{1}^{B},\chi_{2}^{B},\chi_{3}^{B})$ space, and 
    \begin{align}
        \left\langle \boldsymbol{\sigma_{1}}\cdot \boldsymbol{\sigma_{2}} \right\rangle&=
        \begin{bmatrix}
            1 & 0 & 0 & 0 & 0 & 0 \\
            0 & 1 & 0 & 0 & 0 & 0 \\
            0 & 0 & 1 & 0 & 0 & 0 \\
            0 & 0 & 0 & 1 & 0 & 0 \\
            0 & 0 & 0 & 0 & -3 & 0 \\
            0 & 0 & 0 & 0 & 0 & -3
        \end{bmatrix} \\ 
        \left\langle \boldsymbol{\sigma_{1}}\cdot \boldsymbol{\sigma_{3}} \right\rangle&=
        \begin{bmatrix}
            1 & 0 & 0 & 0 & 0 & 0 \\
            0 & -1 & 0 & \sqrt{2} & -\sqrt{2} & 0 \\
            0 & 0 & -2 & 0 & 0 & -\sqrt{3} \\
            0 & \sqrt{2} & 0 & 0 & 1 & 0 \\
            0 & -\sqrt{2} & 0 & 1 & 0 & 0 \\
            0 & 0 & -\sqrt{3} & 0 & 0 & 0
        \end{bmatrix} \\ 
        \left\langle \boldsymbol{\sigma_{1}}\cdot \boldsymbol{\sigma_{4}} \right\rangle&=
        \begin{bmatrix}
            1 & 0 & 0 & 0 & 0 & 0 \\
            0 & -1 & 0 & -\sqrt{2} & -\sqrt{2} & 0 \\
            0 & 0 & -2 & 0 & 0 & \sqrt{3} \\
            0 & -\sqrt{2} & 0 & 0 & -1 & 0 \\
            0 & -\sqrt{2} & 0 & -1 & 0 & 0 \\
            0 & 0 & \sqrt{3} & 0 & 0 & 0
        \end{bmatrix} \\ 
        \left\langle \boldsymbol{\sigma_{2}}\cdot \boldsymbol{\sigma_{3}} \right\rangle&=
        \begin{bmatrix}
            1 & 0 & 0 & 0 & 0 & 0 \\
            0 & -1 & 0 & \sqrt{2} & \sqrt{2} & 0 \\
            0 & 0 & -2 & 0 & 0 & \sqrt{3} \\
            0 & \sqrt{2} & 0 & 0 & -1 & 0 \\
            0 & \sqrt{2} & 0 & -1 & 0 & 0 \\
            0 & 0 & \sqrt{3} & 0 & 0 & 0
        \end{bmatrix} \\ 
        \left\langle \boldsymbol{\sigma_{2}}\cdot \boldsymbol{\sigma_{4}} \right\rangle&=
        \begin{bmatrix}
            1 & 0 & 0 & 0 & 0 & 0 \\
            0 & -1 & 0 & -\sqrt{2} & \sqrt{2} & 0 \\
            0 & 0 & -2 & 0 & 0 & -\sqrt{3} \\
            0 & -\sqrt{2} & 0 & 0 & 1 & 0 \\
            0 & \sqrt{2} & 0 & 1 & 0 & 0 \\
            0 & 0 & -\sqrt{3} & 0 & 0 & 0
        \end{bmatrix} \\ 
        \left\langle \boldsymbol{\sigma_{3}}\cdot \boldsymbol{\sigma_{4}} \right\rangle&=
        \begin{bmatrix}
            1 & 0 & 0 & 0 & 0 & 0 \\
            0 & 1 & 0 & 0 & 0 & 0 \\
            0 & 0 & 1 & 0 & 0 & 0 \\
            0 & 0 & 0 & -3 & 0 & 0 \\
            0 & 0 & 0 & 0 & 1 & 0 \\
            0 & 0 & 0 & 0 & 0 & -3
        \end{bmatrix} \label{sfcT}
    \end{align}
for tetraquark in $(\chi_{1}^{T},\chi_{2}^{T},\chi_{3}^{T},\chi_{4}^{T},\chi_{5}^{T},\chi_{6}^{T})$ space.

For now, the CMI for a hadronic state can be expressed by summing up all the $C_{ij}$ by color and spin factors, 
which are diagonal elements of matrices in Eq.~(\ref{cfcM}-\ref{cfcT}) and Eq.~(\ref{sfcM}-\ref{sfcT}).
The off-diagonal elements are used for chromomagnetic mixing.


\medskip
\section*{Appendix C}\label{AppendixC}
\setcounter{equation}{0}
\renewcommand{\theequation}{C\arabic{equation}}
According to Appendix A and B, we list CMI matrices of all the chromomagnetic mixing involved in this paper, such as
\begin{equation}\label{mixing23}
    \begin{pmatrix}
        \frac{8}{3}C_{12}-\frac{16}{3}C_{13}-\frac{16}{3}C_{23} & -\frac{8\sqrt{3}}{3}C_{13}+\frac{8\sqrt{3}}{3}C_{23} \\
        -\frac{8\sqrt{3}}{3}C_{13}+\frac{8\sqrt{3}}{3}C_{23} & -8C_{12} \\
    \end{pmatrix}
\end{equation}
for $\left(\phi^{B}\chi_{2}^{B},\phi^{B}\chi_{3}^{B}\right)$ mixing,
\begin{equation}\label{21mixing36}
    \begin{pmatrix}
        \frac{8}{3}(\alpha-\beta)   & 2\sqrt{6}\beta \\
        2\sqrt{6}\beta              & 4\alpha
    \end{pmatrix}
\end{equation}
for $\left(\phi_{2}^{T}\chi_{3}^{T},\phi_{1}^{T}\chi_{6}^{T}\right)$ mixing,
\begin{equation}\label{21mixing54}
    \begin{pmatrix}
        -\frac{8}{3}\theta  & -2\sqrt{2}\beta \\
        -2\sqrt{2}\beta     & -\frac{4}{3}\eta
    \end{pmatrix}
\end{equation}
for $\left(\phi_{2}^{T}\chi_{5}^{T},\phi_{1}^{T}\chi_{4}^{T}\right)$ mixing,
\begin{equation}\label{12mixing36}
    \begin{pmatrix}
        -\frac{4}{3}(\alpha+5\beta) & 2\sqrt{6}\beta \\
        2\sqrt{6}\beta              & -8\alpha
    \end{pmatrix}
\end{equation}
for $\left(\phi_{1}^{T}\chi_{3}^{T},\phi_{2}^{T}\chi_{6}^{T}\right)$ mixing,
\begin{equation}\label{221mixing245}
    \begin{pmatrix}
        \frac{4}{3}(2\alpha-\beta)  & \frac{4\sqrt{2}}{3}\delta & 4\delta \\
        \frac{4\sqrt{2}}{3}\delta   & \frac{8}{3}\eta           & -2\sqrt{2}\beta \\
        4\delta                     & -2\sqrt{2}\beta           & \frac{4}{3}\beta
    \end{pmatrix}
\end{equation}
for $\left(\phi_{2}^{T}\chi_{2}^{T},\phi_{2}^{T}\chi_{4}^{T},\phi_{1}^{T}\chi_{5}^{T}\right)$ mixing,
\begin{equation}\label{221mixing254}
    \begin{pmatrix}
        \frac{4}{3}(2\alpha-\beta)  & -\frac{4\sqrt{2}}{3}\gamma    & -4\gamma \\
        -\frac{4\sqrt{2}}{3}\gamma  & -\frac{8}{3}\theta            & -2\sqrt{2}\beta \\
        -4\gamma                    & -2\sqrt{2}\beta               & -\frac{4}{3}\eta
    \end{pmatrix}
\end{equation}
for $\left(\phi_{2}^{T}\chi_{2}^{T},\phi_{2}^{T}\chi_{5}^{T},\phi_{1}^{T}\chi_{1}^{T}\right)$ mixing,
\begin{equation}\label{112mixing245}
    \begin{pmatrix}
        -\frac{2}{3}(2\alpha+5\beta)    & \frac{10\sqrt{2}}{3}\delta    & 4\delta \\
        \frac{10\sqrt{2}}{3}\delta      & -\frac{4}{3}\eta              & -2\sqrt{2}\beta \\
        4\delta                         & -2\sqrt{2}\beta               & -\frac{8}{3}\theta
    \end{pmatrix}
\end{equation}
for $\left(\phi_{1}^{T}\chi_{2}^{T},\phi_{1}^{T}\chi_{4}^{T},\phi_{2}^{T}\chi_{5}^{T}\right)$ mixing,
where $\alpha=C_{12}+C_{34}$, $\beta=C_{13}+C_{14}+C_{23}+C_{24}$, $\gamma=C_{13}+C_{14}-C_{23}-C_{24}$, $\delta=C_{13}-C_{14}+C_{23}-C_{24}$, 
$\eta=C_{12}-3C_{34}$ and $\theta=3C_{12}-C_{34}$, refering Ref.[].

And we also present binding energies
\begin{equation}
    B_{12}+B_{13}+B_{23}
\end{equation}
for $\phi^{B}$,
\begin{equation}
    -\frac{1}{2}B_{12}+\frac{5}{4}B_{13}+\frac{5}{4}B_{14}+\frac{5}{4}B_{23}+\frac{5}{4}B_{24}-\frac{1}{2}B_{34}
\end{equation}
for $\phi_{1}^{T}$, and
\begin{equation}
    B_{12}+\frac{1}{2}B_{13}+\frac{1}{2}B_{14}+\frac{1}{2}B_{23}+\frac{1}{2}B_{24}+B_{34}
\end{equation}
for $\phi_{2}^{T}$ respectively, where $B_{ij}$ takes value from Eq.~(\ref{bingdings}) and equals 0 if there is no short-disctance binding.


\begin{thebibliography}{99}
    \bibliographystyle{unsrt}

%\bibitem{Tanabashi:D18} P.A.~Zyla \textit{et al.} [Particle Data Group], Review of Particle Physics, Prog. Theor. Exp. Phys. 2020, 083C01 (2020)

\end{thebibliography}

\end{document}
